\documentclass{article} % For LaTeX2e
\usepackage{iclr2027_conference,times}

% Optional math commands from https://github.com/goodfeli/dlbook_notation.
\input{math_commands.tex}

\usepackage{hyperref}
\usepackage{url}

% Packages required for the paper's own content (tables, math notation, figures).
\usepackage{amsmath}
\usepackage{amssymb}
\usepackage{booktabs}
\usepackage{multirow}
\usepackage{graphicx}

\newcommand{\prism}{\textsc{Prism}}
\newcommand{\tbd}{--}

\title{Disentangling the Spectrum with PRISM: A Dual-Stream Physics-Informed VAE for Decoupled Representational Learning of Hyperspectral Images}

\author{Anonymous authors \\ Paper under double-blind review}

%\iclrfinalcopy % Uncomment for camera-ready version, but NOT for submission.
\begin{document}

\maketitle

\begin{abstract}
Hyperspectral imaging (HSI) captures dense reflectance spectra across hundreds of contiguous spectral bands. Robust representation learning of this data is central to applications like remote sensing, mineral mapping, atmospheric characterization etc. However sensor degradation, transmission failures, and external noise routinely corrupt these cubes, motivating a need for variational representation learning, typically via VAEs, to compress and probabilistically reconstruct corrupted spectra. Conventional VAEs apply convolutions that mix spatial and spectral information indiscriminately, blurring spectral signatures and causing posterior collapse or physically improbable hallucinations. Thus, we introduce PRISM (Physics-Informed Representation for Isolated Spectral-Spatial Modeling). PRISM is a dual-stream VAE routing spatial and spectral information through two independently-supervised spatial and spectral branches that re-couple only under a Spectral Angle Mapper (SAM) bounded fusion. This lets PRISM retain the spatial fidelity of standard 2D VAEs without sacrificing chemical accuracy; gaining 6.949 dB PSNR and 0.112 degree SAM improvement over the strongest rival architectures across the ablation grid. This decoupled latent is designed as a backbone for Latent Diffusion Models (LDMs) for purifying deep-space transmission corruptions and generating omitted frames. It already supports chemical interpolation for mineral discovery while offering a scalable framework for other physics-heavy domains like precision agriculture and climatology.
\end{abstract}

\section{Introduction}

Hyperspectral images (HSI) are imaging spectrometer products that record, at every pixel, a nearly-continuous reflectance spectrum sampled across hundreds of narrow, contiguous wavelength bands rather than the three broadband channels of an ordinary photograph. Where an RGB camera integrates incident light over wide red, green and blue windows, an HSI sensor resolves the fine spectral structure -- absorption dips, shoulders and continuum slopes -- that is diagnostic of the physical and chemical composition of the imaged surface. A single HSI cube is therefore simultaneously a spatial image and, at every pixel, a spectroscopic measurement.

This dual nature is what makes HSI valuable across a wide range of remote-sensing domains. On Earth, HSI underlies precision agriculture (crop stress and nutrient mapping), environmental and atmospheric monitoring, and mineral exploration. Off Earth, imaging spectrometers such as the Chandrayaan-2 Imaging Infra-Red Spectrometer (IIRS) and the Compact Reconnaissance Imaging Spectrometer for Mars (CRISM) are the primary instruments used to map lunar and Martian surface mineralogy and hydration state from orbit \citep{chowdhury2020iirs, murchie2007crism}, while airborne sensors such as AVIRIS \citep{green1998aviris} connect these planetary use cases to a mature terrestrial calibration and validation literature.

Because an HSI cube is large, redundant along the spectral axis, and expensive to transmit or store in full, representation learning has become central to working with it: compressing a cube into a compact latent code, learning embeddings for downstream classification or unmixing, and -- increasingly -- producing the latent inputs that modeling frameworks such as Latent Diffusion Models (LDMs) require as their starting point rather than operating on full-resolution pixels directly.

A purely deterministic compressor, however, is a poor fit for a further recurring problem in HSI: sensor degradation, transmission dropouts, and instrument or environmental noise routinely corrupt real cubes, and a useful representation should not just compress a clean cube but plausibly \emph{fix or regenerate} the corrupted or missing parts of one. This is precisely the objective that variational representation learning is built for: a Variational Autoencoder (VAE) \citep{kingma2014vae} learns a smooth, continuous, and regularized latent space rather than memorizing discrete codes, which is exactly what lets it interpolate between and probabilistically regenerate spectra it has not seen verbatim -- denoising a corrupted patch, filling a dropped band, or synthesizing a chemically plausible signature between two known materials -- rather than only reproducing what a deterministic autoencoder was shown at training time.

Realizing that promise in practice, however, runs into a structural problem with how standard VAEs are built for imagery. The dominant convolutional VAE treats an HSI cube exactly as it would treat an RGB image: 2D convolutions slide a kernel over the spatial $(H, W)$ plane and, when applied naively to a cube with hundreds of bands, mix spatial neighbourhoods and spectral channels indiscriminately in the same operation. This is an excellent inductive bias for spatial denoising and structure, but it is the wrong one for spectral fidelity: it blurs the very absorption features that carry chemical meaning, and in the aggressive bottlenecks a VAE latent requires it manifests as posterior collapse or as reconstructions whose reconstructed spectra are spatially plausible but chemically wrong -- a physically improbable hallucination that is far more consequential in a scientific pipeline than an ordinary image artifact would be.

We introduce PRISM (Physics-Informed Representation for Isolated Spectral-Spatial Modeling) to resolve this conflict directly rather than trade it off. PRISM is a dual-stream VAE that routes spatial and per-pixel spectral information through two independently-supervised encoder-decoder branches, and re-couples their reconstructions only at a single, physics-bounded fusion step governed by a differentiable Spectral Angle Mapper (SAM) penalty -- so that a spatial denoiser and a spectral-fidelity model are each free to specialize, and are only ever combined under a term that directly penalizes the chemical hallucination a naive fusion would otherwise permit.

Our contributions are:
\begin{enumerate}
\item A split-stream architecture for self-supervised representation learning of hyperspectral cubes, in which spatial context and per-pixel spectral structure are encoded, reparameterized, and decoded independently before being re-coupled at a single fusion step.
\item A Spectral-Angle-Mapper-bounded physics-informed fusion and objective, which penalizes chemically-hallucinated reconstructions directly rather than relying on reconstruction error alone.
\item A falsification-oriented ablation protocol -- latent-rate matching, posterior-collapse detection, and a spatial-reliance shuffle test -- together with a downstream-readiness benchmark suite (latent noise-injection robustness, chemical-interpolation smoothness, and spectral band-masking recovery) that evaluates PRISM's latent as a candidate LDM backbone rather than on reconstruction fidelity alone.
\end{enumerate}

\section{Related Works}

\paragraph{Representation learning and spectral unmixing for HSI.} A large body of work applies deep representation learning to hyperspectral data for compression, classification, and spectral unmixing without committing to a VAE-specific formulation. Convolutional architectures adapted from natural-image classification learn joint spatial-spectral features directly from raw cubes \citep{chen2016deep}, and hybrid spatial-spectral CNNs such as HybridSN combine 3D and 2D convolutions in a single feature hierarchy to trade off the two kinds of structure \citep{roy2020hybridsn}. Transformer-based encoders have more recently been proposed to model long-range spectral dependencies directly, as in SpectralFormer \citep{hong2022spectralformer}, and spatial-spectral recurrent transformers have been used specifically for HSI denoising and restoration \citep{fu2024sst}, alongside earlier convolutional restoration networks such as HSI-DeNet \citep{chang2019hsidenet}. On the unmixing side, autoencoder-based unmixing recovers a small set of endmember spectra and per-pixel abundances by construction, an objective closely related to representation learning even though it does not train a generative latent space \citep{su2019daen, palsson2018unmixing}. The Spectral Angle Mapper itself, which PRISM adopts as a fusion-bounding loss, originates in this classical spectral-matching literature \citep{kruse1993sips} rather than in any VAE work, underscoring that it measures a property -- spectral direction -- that is orthogonal to whatever architecture is used to produce a reconstruction.

\paragraph{Differentiable physics-penalty training.} A separate line of work constrains network outputs with a differentiable physical prior rather than relying on architecture or capacity alone to enforce physical plausibility. Physics-informed neural networks add a residual of a governing differential equation directly into the training loss so that the network's output is penalized whenever it violates known physics \citep{raissi2019pinn}. Closer to our domain, physics-guided autoencoders for planetary spectroscopy have used a radiative-transfer model (the Hapke model) as a differentiable penalty during unmixing of lunar hyperspectral data \citep{lin2026hapke}, which is structurally the same idea as PRISM's SAM-bounded fusion: a physically-motivated, differentiable term added to an otherwise standard reconstruction objective, rather than a hand-engineered post-processing step.

\paragraph{Specialized VAE architectures for HSI.} Most directly relevant are VAE variants adapted specifically to the geometry of hyperspectral data, which is also where our three baseline architectures are anchored. The foundational VAE and $\beta$-VAE formulations \citep{kingma2014vae, higgins2017betavae} motivate treating the latent bottleneck itself, not just the decoder capacity, as the primary object of control -- the same principle behind this paper's latent-rate-matched ablation protocol (\S\ref{sec:datasets}). Stable Diffusion's \texttt{AutoencoderKL} \citep{rombach2022ldm} is the reference architecture for our 2D spatial baseline: a purely 2D-convolutional VAE that treats the input as a thick image and is exactly the design whose spatial-spectral mixing motivates PRISM's split. At the opposite extreme, per-pixel spectral (V)AEs process each pixel's spectrum independently through a shared MLP, discarding spatial structure entirely in exchange for chemically pristine per-pixel codes; \citet{palsson2018unmixing} train such an autoencoder for hyperspectral unmixing with an objective that directly rewards reconstruction and abundance fidelity over raw pixel error, and \citet{liu2022dualfreq} adopt the same four-fully-connected-layer per-pixel design for anomaly detection in transformed hyperspectral imagery -- together anchoring our 1D pixelwise baseline. In between, fully-3D convolutional autoencoders apply 3D kernels across the joint spatial-spectral volume to learn features that span both axes simultaneously \citep{chen2016deep}; our 3D spatio-spectral baseline follows this design, at the parameter and compute cost of a kernel that structurally averages neighbouring bands together with neighbouring pixels. PRISM differs from all three by refusing to let any single kernel see both axes at once: the spatial and spectral streams are architecturally isolated until a single, physics-bounded fusion step recombines them.

\section{Methodology}

\subsection{PRISM Architecture}
\label{sec:architecture}

PRISM is a dual-stream variational autoencoder. Given an input patch $\mathbf{x} \in \mathbb{R}^{H \times W \times C}$ ($H=W=64$; $C$ the sensor's band count), a \emph{spatial stream} and a \emph{spectral stream} each independently encode, reparameterize, and decode a reconstruction of the full patch; the two reconstructions are combined only at a final, physics-bounded fusion step. Figure~\ref{fig:l1dfd} (level-1 data-flow diagram) and Figure~\ref{fig:l2dfd} (level-2, tensor-shape-annotated) summarize the architecture; both are supplied by the authors and are placeholders in this draft.

\begin{figure}[t]
\centering
\fbox{\parbox{0.85\linewidth}{\centering\vspace{2.2cm}\textit{Level-1 architecture diagram (spatial stream / spectral stream / shared reparameterization / physics-bounded fusion) -- to be inserted.}\vspace{2.2cm}}}
\caption{PRISM architecture overview (level-1 data-flow diagram).}
\label{fig:l1dfd}
\end{figure}

\paragraph{Spatial stream.} The spatial stream is responsible for spatial context and denoising. Its encoder first projects every pixel's $C$-band spectrum onto a $r{=}64$-channel local representation with a $3{\times}3$ 2D convolution applied across all bands at once ($\mathrm{Conv2d}(C{\to}r, k{=}3, p{=}1)$) -- a design we refer to as PRISM's \emph{non-linear} (2D) spatial projection, in contrast to a per-pixel $1{\times}1$-receptive-field alternative that projects each pixel's spectrum independently of its neighbours. Three stride-2 convolutional blocks ($\mathrm{Conv2d}(k{=}3,s{=}2,p{=}1)$ + ReLU, channels $r{\to}2r{\to}4r{\to}8r$) then downsample the spatial grid $64{\to}32{\to}16{\to}8$, and a final $1{\times}1$ convolution projects to a $(B, 2d_s, 8, 8)$ mean/log-variance map, where $d_s$ is the spatial-stream's per-dataset latent width. The decoder mirrors this exactly: a $1{\times}1$ convolution back to $8r$ channels, three stride-2 transposed-convolution blocks ($\mathrm{ConvTranspose2d}(k{=}4,s{=}2,p{=}1)$ + ReLU) upsampling $8{\to}16{\to}32{\to}64$, and a final $3{\times}3$ convolution back to $C$ bands followed by a sigmoid. The spatial stream's latent is therefore an $8{\times}8$ spatial grid of $d_s$-dimensional codes, $\mathbf{z}_s \in \mathbb{R}^{d_s \times 8 \times 8}$ -- an addressable spatial map, not a single global vector, so that the spatial stream keeps the ability to vary its code across the patch.

\paragraph{Spectral stream.} The spectral stream is responsible for per-pixel chemical fidelity and never mixes information across pixels. Every pixel's spectrum is treated independently: the encoder applies two stride-2 1D convolutions along the spectral axis ($\mathrm{Conv1d}(k{=}4,s{=}2,p{=}1)$ + ReLU, channels $1{\to}32{\to}64$), flattens, and projects through a linear layer to a $2d_p$-dimensional mean/log-variance vector per pixel, where $d_p$ is the spectral-stream's per-dataset latent width. Reshaped back onto the spatial grid, this gives a per-pixel latent map $\mathbf{z}_p \in \mathbb{R}^{d_p \times H \times W}$. The decoder mirrors the encoder with two stride-2 1D transposed convolutions, ending in a sigmoid. Because every layer in this stream operates on one pixel's $C$-vector at a time, the spectral stream is architecturally incapable of blurring one pixel's spectrum into another's.

\paragraph{Shared reparameterization.} Both streams' encoders emit a $(mean, \log\text{-}variance)$ pair, chunked along the channel axis; a single shared reparameterization function samples $\mathbf{z} = \boldsymbol{\mu} + \boldsymbol{\sigma} \odot \boldsymbol{\varepsilon}$, $\boldsymbol{\varepsilon} \sim \mathcal{N}(0, I)$, with $\log\text{-}variance$ clamped to $[-30, 20]$ for numerical stability, applied identically to $\mathbf{z}_s$ and $\mathbf{z}_p$.

\begin{figure}[t]
\centering
\fbox{\parbox{0.85\linewidth}{\centering\vspace{2.2cm}\textit{Level-2 tensor-flow diagram (per-layer shapes through both streams and the fusion) -- to be inserted.}\vspace{2.2cm}}}
\caption{PRISM tensor-shape flow (level-2 data-flow diagram).}
\label{fig:l2dfd}
\end{figure}

\paragraph{Physics-bounded fusion.} The two streams' decoded reconstructions $\hat{\mathbf{x}}_s$ (spatial) and $\hat{\mathbf{x}}_p$ (spectral), each already in $[0,1]$ from its own sigmoid, are combined by a spatially-adaptive gate: a two-layer $3{\times}3$-convolutional network computes a per-pixel, per-band gate $\boldsymbol{\alpha} = \sigma(g([\hat{\mathbf{x}}_s; \hat{\mathbf{x}}_p]))$ from the concatenated reconstructions, and the fused output is the convex combination
\begin{equation}
\hat{\mathbf{x}}_f = \boldsymbol{\alpha} \odot \hat{\mathbf{x}}_s + (1 - \boldsymbol{\alpha}) \odot \hat{\mathbf{x}}_p .
\end{equation}
Because $\boldsymbol{\alpha} \in [0,1]$ and both branch reconstructions are already bounded, $\hat{\mathbf{x}}_f$ requires no further activation. The gate is trained jointly with the rest of the network and is bounded by the SAM term described next, so the fusion cannot learn to blend the two streams in a way that produces a spectrum pointing in the wrong chemical direction, even if doing so would reduce pixel-wise error.

\subsection{Objective Function and Physics-Informed Fusion}
\label{sec:objective}

PRISM's training objective adds a spectral-fidelity penalty to a standard multi-term ELBO. Writing $\hat{\mathbf{x}}_f$, $\hat{\mathbf{x}}_s$, $\hat{\mathbf{x}}_p$ for the fused, spatial-only, and spectral-only reconstructions,
\begin{equation}
\mathcal{L}_{\mathrm{rec}} = \frac{0.5 \cdot \mathrm{MSE}(\hat{\mathbf{x}}_f, \mathbf{x}) + w \cdot \mathrm{MSE}(\hat{\mathbf{x}}_s, \mathbf{x}) + w \cdot \mathrm{MSE}(\hat{\mathbf{x}}_p, \mathbf{x})}{0.5 + 2w}, \qquad w = 0.1,
\end{equation}
a normalized weighted mean (weights sum to 1) with an effective $5{:}1{:}1$ ratio between the fused term and each branch's own auxiliary reconstruction term. The auxiliary terms exist so that each stream is individually supervised -- neither branch can rely entirely on the other to reconstruct the patch -- while the normalization keeps $\mathcal{L}_{\mathrm{rec}}$'s scale directly comparable to a single-stream baseline's plain MSE, so that a shared $\beta$ and $\lambda_{\mathrm{phys}}$ (below) mean the same thing across every architecture in the ablation. The KL term is likewise a mean over both streams' divergences,
\begin{equation}
\mathcal{L}_{\mathrm{KL}} = \tfrac{1}{2}\left(\mathrm{KL}(\mathbf{z}_s) + \mathrm{KL}(\mathbf{z}_p)\right), \qquad \mathrm{KL}(\mathbf{z}) = -\tfrac{1}{2}\,\mathrm{mean}\left(1 + \log\sigma^2 - \mu^2 - \sigma^2\right),
\end{equation}
reduced over every element of the latent so that latent size does not bias its magnitude across streams of very different shape. The physics prior is the Spectral Angle Mapper, applied only to the \emph{fused} reconstruction:
\begin{equation}
\mathcal{L}_{\mathrm{SAM}}(\mathbf{x}, \hat{\mathbf{x}}_f) = \mathrm{mean}_{i,j}\; \arccos\!\left(\frac{\langle \mathbf{x}_{ij}, \hat{\mathbf{x}}_{f,ij}\rangle}{\left(\lVert\mathbf{x}_{ij}\rVert + \epsilon\right)\left(\lVert\hat{\mathbf{x}}_{f,ij}\rVert + \epsilon\right) + \epsilon}\right), \qquad \epsilon = 10^{-8},
\end{equation}
with the cosine additionally clamped to $[-1+\epsilon, 1-\epsilon]$ before the arccosine. SAM is invariant to illumination scaling, so it penalizes a reconstruction whose spectral \emph{direction} differs from the input's regardless of its magnitude -- exactly the chemical-hallucination failure mode a naive fusion would otherwise be free to produce. (A pixel whose true spectral energy falls below $\epsilon$ contributes exactly $\pi/2$ to $\mathcal{L}_{\mathrm{SAM}}$ irrespective of the model's prediction; we report both the raw SAM and an energy-masked \emph{SAM-valid} that excludes such pixels, and use SAM-valid for every cross-model and cross-dataset comparison in \S\ref{sec:results}.) The full objective is
\begin{equation}
\mathcal{L} = \mathcal{L}_{\mathrm{rec}} + \beta\,\mathcal{L}_{\mathrm{KL}} + \lambda_{\mathrm{phys}}\,\mathcal{L}_{\mathrm{SAM}},
\end{equation}
with $\beta = 10^{-3}$ and $\lambda_{\mathrm{phys}} = 0.3$ held constant throughout training (no annealing) and identical across every model and dataset in the ablation grid.

\subsection{Experiment Setup}
\label{sec:datasets}

\paragraph{Ablation grid.} We compare PRISM against three baseline architectures on three hyperspectral datasets, all trained under an identical pipeline. \textbf{2D Spatial VAE} is a direct adaptation of Stable Diffusion's \texttt{AutoencoderKL} \citep{rombach2022ldm}: purely 2D convolutions treat the cube as a thick image, downsampling and reconstructing the $(H,W)$ grid with no explicit spectral-axis structure. \textbf{3D Spatio-Spectral VAE} follows the fully-3D convolutional autoencoder line for HSI \citep{chen2016deep}: the patch is lifted to a single-channel $(1, C, H, W)$ volume and processed entirely with 3D convolutions and transposed convolutions, so every kernel spans both spatial and spectral axes at once. \textbf{1D Pixelwise VAE} follows the per-pixel spectral autoencoder line \citep{palsson2018unmixing, liu2022dualfreq}: every $(H,W)$ location is folded into the batch dimension and each pixel's spectrum is passed independently through a shared fully-connected encoder/decoder (four layers, LayerNorm + GELU), so the model never sees the 2D grid at all. All three architectures were initialized at the widths reported in their respective papers rather than being parameter-matched to PRISM: the 2D Spatial VAE uses \texttt{AutoencoderKL}'s base channel width of 128 over three downsampling stages; the 3D Spatio-Spectral VAE uses a compact base width of 24 filters, representative of the published 3D convolutional autoencoder line for HSI classification (the exact per-layer filter counts were not recoverable from that architecture's original description, so this width is a representative rather than an exact reproduction); and the 1D Pixelwise VAE uses hidden widths $(512, 256, 128)$, matching the four-fully-connected-layer design used for spectral unmixing and anomaly detection in the cited works. All five architectures (three baselines, each trainable under a standard ELBO objective or the physics-augmented objective of \S\ref{sec:objective}, plus PRISM, for which the SAM prior is intrinsic and not ablated) implement one shared model-agnostic contract (\texttt{forward}, \texttt{loss\_terms}, \texttt{reconstruct}, \texttt{encode\_latents}, \texttt{decode\_latents}), so the training and evaluation code never branches on model identity. Every architecture is trained at two random seeds (67, 69) under the physics-augmented objective; the three baselines are additionally trained once under a standard-ELBO objective, to isolate the effect of the SAM term itself.

Latent geometry differs by design across the four architectures -- that difference \emph{is} what the ablation tests -- but latent \emph{rate} is matched: for each dataset, every model's bottleneck is solved to land on the same total latent-element budget $T$ (the nearest multiple of $H \cdot W = 4096$ to a $64{:}1$ compression target), so that no architecture is being compared at a systematically looser or tighter information bottleneck than another. On IIRS, every model lands on $T = 16{,}384$ elements exactly ($64.0{:}1$); on AVIRIS and CRISM, every model lands on $T = 28{,}672$ ($60.6{:}1$ and $65.1{:}1$ respectively, both within the preregistered $\pm25\%$ tolerance of the $64{:}1$ target). Parameter count is reported alongside every result rather than matched: PRISM has $5.20$--$5.95$M parameters depending on dataset, the 2D Spatial VAE $23.4$--$24.5$M, the 3D Spatio-Spectral VAE a constant $3.10$M, and the 1D Pixelwise VAE $0.60$--$0.80$M (Table~\ref{tab:table1}).

\paragraph{Data pipeline.} We split each hyperspectral cube into smaller $64{\times}64{\times}C$ cubes (patches), extracted at a stride of 48 pixels (25\% overlap) from region-disjoint 70/15/15 train/validation/test slices of each source scene, so that a training patch and a test patch are never drawn from adjacent, overlapping terrain. We then create a combined packed set of a fixed number of patches per dataset to train on: the training split of every dataset is capped at 7{,}000 patches (or its full size where smaller), subsampled proportionally across source scenes with a fixed seed, so that no dataset dominates training simply by having more raw scenes available; validation and test splits are never capped. All models are trained for 60 epochs (AdamW, learning rate $10^{-3}$, 5-epoch linear warm-up followed by cosine annealing, gradient-norm clipped at 1.0) with early stopping at a patience of 3 epochs on either of two monitored criteria (validation SAM or validation reconstruction MSE); every cell reports the checkpoint selected on validation SAM, evaluated once at each of the two training seeds.

\begin{table}[t]
\centering
\small
\begin{tabular}{lcccccc}
\toprule
\textbf{Dataset} & \textbf{Instrument} & \textbf{Bands} & \textbf{Wavelength} & \textbf{Train} & \textbf{Valid} & \textbf{Test} \\
\midrule
IIRS & Chandrayaan-2 IIRS \citep{chowdhury2020iirs} & 256 & $\sim$0.8--5.0 $\mu$m & 14{,}624 & 3{,}084 & 3{,}084 \\
AVIRIS & Airborne AVIRIS \citep{green1998aviris} & 424 & 0.4--2.5 $\mu$m & 7{,}000 & 2{,}180 & 1{,}950 \\
CRISM & MRO CRISM \citep{murchie2007crism} & 456 & 0.362--3.92 $\mu$m & 2{,}561 & 369 & 369 \\
\bottomrule
\end{tabular}
\caption{Datasets used in the ablation grid, after preprocessing and patch extraction. AVIRIS's training split is capped to the 7{,}000-patch budget shared by the larger datasets; CRISM's smaller scene count is not capped and uses its full 2{,}561-patch training split. Band counts are the number of contiguous reflective channels retained by the packing pipeline (after removing non-reflective or degenerate detector channels), not the raw sensor channel count.}
\label{tab:datasets}
\end{table}

\section{Results and Discussion}
\label{sec:results}

\paragraph{Mechanism diagnostics.} Before reading any reconstruction number, we first check that every trained cell is measuring a real, non-degenerate model rather than an artifact of the training or evaluation pipeline, using three diagnostics run on every frozen checkpoint. \emph{Latent-rate audit} confirms every cell lands within the preregistered $\pm25\%$ tolerance of its dataset's common latent budget $T$ (\S\ref{sec:datasets}); a rate-unmatched cell cannot be attributed a win or loss on architecture alone, since a looser bottleneck trivially helps reconstruction. \emph{Posterior-collapse detection} checks that a model's latent is actually used, via the fraction of latent units carrying non-trivial KL divergence and a latent-swap test (decoding one patch's reconstruction from a different patch's code): a collapsed decoder emits a near-constant output regardless of its input latent, and such cells are excluded from every comparison rather than silently scored. \emph{Spatial-reliance shuffle} permutes the spatial grid of pixels (each pixel's own spectrum kept intact) and re-scores the shuffled reconstruction against the shuffled input, reporting the resulting relative SAM degradation (SRI): a model with no genuine use of spatial context is exactly unaffected by the shuffle, which is why the 1D Pixelwise VAE serves as this diagnostic's positive control (its intact and shuffled scores must match to within numerical tolerance, or the diagnostic itself, not the model, is at fault). Two supporting diagnostics complete the picture and are reported in full in Appendix~\ref{app:diagnostics}: a tiered trivial-predictor floor (does a trained model actually beat a data-only predictor that never sees the training set, at the same evaluation cost) and a spectral band-masking recovery test that anticipates Table~\ref{tab:table4} below. None of these diagnostics carries a pass/fail verdict in this draft; each is reported as a number precisely because a hard threshold on this kind of check has previously proven miscalibrated on smooth, highly-redundant hyperspectral data (a model can lose to its own patch-level spatial mean at a tight compression ratio without that being a fault in the model), and we report the diagnostic honestly rather than gate results on a threshold we do not yet trust.

This section reports four tables: reconstruction quality (Table~\ref{tab:table1}), recovery from injected latent noise (Table~\ref{tab:table2}), chemical-interpolation smoothness (Table~\ref{tab:table3}), and spectral band-masking recovery (Table~\ref{tab:table4}). Every table is seed-averaged over both training seeds (67, 69), reads the physics-objective, validation-SAM-selected checkpoint for every model, and covers IIRS, AVIRIS, and CRISM.

\begin{table}[t]
\centering
\small
\setlength{\tabcolsep}{4.5pt}
\resizebox{\linewidth}{!}{%
\begin{tabular}{llccccccc}
\toprule
\textbf{Dataset} & \textbf{Model} & \textbf{SAM-valid ($^\circ$)$\downarrow$} & \textbf{PSNR$\uparrow$} & \textbf{SSIM$\uparrow$} & \textbf{SID$\downarrow$} & \textbf{SCC$\uparrow$} & \textbf{Q2$^n$$\uparrow$} & \textbf{Params (M)} \\
\midrule
% GENERATED by paper/iclr/v2-claude/make_tables.py — DO NOT EDIT BY HAND
% source: /home/minikazuki/Documents/Research Papers/A Dual Stream Physics Informed VAE for hallucination free LDM purification of HSI/results/ablation_table.csv + /home/minikazuki/Documents/Research Papers/A Dual Stream Physics Informed VAE for hallucination free LDM purification of HSI/results/model_params.json | seeds=[67, 69] loss=physics select=sam
\multirow{4}{*}{IIRS} & 2D Spatial VAE & 9.85 & 39.93 & 0.946 & 0.1653 & 0.074 & \textbf{0.708} & 23.40 \\
 & 1D Pixelwise VAE & \textbf{6.93} & 32.47 & 0.825 & \textbf{0.0948} & \textbf{0.222} & 0.639 & 0.60 \\
 & 3D Spatio-Spectral VAE & 11.53 & 39.19 & \textbf{0.955} & 0.1973 & 0.105 & 0.569 & 3.10 \\
 & \textbf{PRISM} & 7.01 & \textbf{39.97} & 0.932 & \tbd{} & \tbd{} & \tbd{} & 5.20 \\
\midrule
\multirow{4}{*}{AVIRIS} & 2D Spatial VAE & 5.43 & 14.78 & 0.720 & 0.0224 & 0.000 & 0.124 & 24.38 \\
 & 1D Pixelwise VAE & \textbf{1.62} & 34.93 & 0.849 & \textbf{0.0017} & \textbf{0.460} & \textbf{0.808} & 0.77 \\
 & 3D Spatio-Spectral VAE & 1.97 & 33.24 & 0.839 & 0.0029 & 0.089 & 0.676 & 3.10 \\
 & \textbf{PRISM} & 1.63 & \textbf{35.83} & \textbf{0.872} & \tbd{} & \tbd{} & \tbd{} & 5.85 \\
\midrule
\multirow{4}{*}{CRISM} & 2D Spatial VAE & 6.12 & 23.77 & 0.801 & 0.3769 & 0.028 & 0.213 & 24.45 \\
 & 1D Pixelwise VAE & 3.71 & 29.46 & 0.760 & 0.1976 & \textbf{0.394} & 0.339 & 0.80 \\
 & 3D Spatio-Spectral VAE & 2.82 & 27.13 & \textbf{0.927} & \textbf{0.0672} & 0.046 & \textbf{0.362} & 3.10 \\
 & \textbf{PRISM} & \textbf{2.65} & \textbf{33.95} & 0.814 & \tbd{} & \tbd{} & \tbd{} & 5.95 \\
\bottomrule

\end{tabular}%
}
\caption{Reconstruction quality on the held-out test split. Best value per dataset and column in bold. SID (spectral information divergence) and SCC (spatial correlation coefficient) are reported alongside PSNR/SSIM/SAM to give a fuller spatial-and-spectral fidelity picture; Q2$^n$ is a combined hypercomplex quality index. Both SID/SCC/Q2$^n$ and Q2$^n$'s comparability depend on each sensor's own band count and are meaningful within a dataset, not across datasets. Cells marked ``\tbd{}'' await a pending re-evaluation pass that had not completed at submission time.}
\label{tab:table1}
\end{table}

Table~\ref{tab:table1} shows PRISM winning PSNR on every one of the three datasets and SAM-valid on two of three, while never losing SSIM by more than a few thousandths on the one dataset (IIRS) where the 3D Spatio-Spectral VAE's SSIM is marginally higher. The pattern the architecture predicts is visible in the baselines: the 2D Spatial VAE is the best or near-best on SSIM everywhere but has by far the worst SAM on every dataset (its SAM-valid is $2$--$5{\times}$ PRISM's), consistent with 2D convolutions blurring spectral direction across pixels even as they preserve spatial structure; the 1D Pixelwise VAE is competitive on SAM on AVIRIS and CRISM but loses badly on IIRS, where it has no spatial context to denoise a harder sensor; and the 3D Spatio-Spectral VAE, at roughly $3$--$5{\times}$ PRISM's parameter count, sits in between on most metrics without dominating any of them. No single baseline is simultaneously competitive on spectral and spatial fidelity across all three sensors; PRISM is, at $5.2$--$6.0$M parameters -- smaller than the 2D Spatial VAE and the 3D Spatio-Spectral VAE, and comparable in order of magnitude to (though larger than) the 1D Pixelwise VAE.

\begin{table}[t]
\centering
\small
\begin{tabular}{llccc}
\toprule
\textbf{Dataset} & \textbf{Model} & \textbf{PSNR@$\sigma{=}0.5\uparrow$} & \textbf{SAM@$\sigma{=}0.5$ ($^\circ$)$\downarrow$} & \textbf{PSNR drop} \\
\midrule
% GENERATED by paper/iclr/v2-claude/make_tables.py — DO NOT EDIT BY HAND
% source: /home/minikazuki/Documents/Research Papers/A Dual Stream Physics Informed VAE for hallucination free LDM purification of HSI/results/downstream_table.csv | seeds=[67, 69] loss=physics select=sam | sigma=0.5 only (psnr_clean/psnr_mid=sigma0/sigma0.5)
\multirow{4}{*}{IIRS} & 2D Spatial VAE & \textbf{36.59} & \textbf{7.91} & 2.40 \\
 & 1D Pixelwise VAE & 29.26 & 9.77 & 1.32 \\
 & 3D Spatio-Spectral VAE & 36.34 & 10.06 & 1.73 \\
 & \textbf{PRISM} & 30.67 & 14.49 & 8.57 \\
\midrule
\multirow{4}{*}{AVIRIS} & 2D Spatial VAE & 15.76 & 4.61 & -0.78 \\
 & 1D Pixelwise VAE & \textbf{26.57} & \textbf{2.20} & 8.23 \\
 & 3D Spatio-Spectral VAE & 26.54 & 4.19 & 8.52 \\
 & \textbf{PRISM} & 20.84 & 4.83 & 14.35 \\
\midrule
\multirow{4}{*}{CRISM} & 2D Spatial VAE & 23.83 & 19.55 & -1.56 \\
 & 1D Pixelwise VAE & 23.09 & \textbf{18.20} & 2.91 \\
 & 3D Spatio-Spectral VAE & \textbf{27.33} & 18.86 & 6.27 \\
 & \textbf{PRISM} & 18.49 & 20.85 & 16.94 \\
\bottomrule

\end{tabular}
\caption{Recovery from Gaussian latent noise injected at $\sigma{=}0.5$ (absolute, added directly to the encoded latent). PSNR drop is measured relative to each model's own clean-latent reconstruction. Lower drop is not unambiguously better in isolation: a model whose clean reconstruction is already poor has little quality left to lose, so PSNR drop must be read together with the absolute PSNR-at-$\sigma{=}0.5$ column, not on its own.}
\label{tab:table2}
\end{table}

Table~\ref{tab:table2} is the one place PRISM's headline advantage does not carry over cleanly, and we report it plainly rather than soften it. PRISM shows the largest raw PSNR drop under $\sigma{=}0.5$ latent noise on all three datasets. Read in isolation this looks like a fragility finding, but the raw drop is confounded by starting quality: the 2D Spatial VAE shows a small or even \emph{negative} drop on AVIRIS and CRISM precisely because its clean reconstruction is already collapsed or near-collapsed in the physics-objective regime (Appendix~\ref{app:diagnostics} identifies these as posterior-collapse cells), so there is little quality left for noise to destroy. Read on absolute terms instead, PRISM's PSNR-at-$\sigma{=}0.5$ remains competitive with, though not always the best among, the baselines that have not collapsed. We take this honestly as a genuine finding about PRISM's decoder rather than an artifact: a high-fidelity, chemically-precise decoder can be more sensitive to a latent-space perturbation of a fixed absolute size than a coarser one is, simply because it has more reconstructed detail available to disturb. We recover and report the full noise-level sweep ($\sigma \in \{0, 0.1, 0.5, 1.0\}$) in a revision once the corresponding intermediate results are regenerated; the present draft reports the single $\sigma{=}0.5$ point specified by our protocol.

\begin{table}[t]
\centering
\small
\begin{tabular}{llcc}
\toprule
\textbf{Dataset} & \textbf{Model} & \textbf{Jaggedness$\downarrow$} & \textbf{Path length} \\
\midrule
% GENERATED by paper/iclr/v2-claude/make_tables.py — DO NOT EDIT BY HAND
% source: /home/minikazuki/Documents/Research Papers/A Dual Stream Physics Informed VAE for hallucination free LDM purification of HSI/results/downstream_table.csv | seeds=[67, 69] loss=physics select=sam
\multirow{4}{*}{IIRS} & 2D Spatial VAE & 0.0095 & 0.488 \\
 & 1D Pixelwise VAE & 0.0243 & 0.381 \\
 & 3D Spatio-Spectral VAE & 0.0053 & 0.548 \\
 & \textbf{PRISM} & \textbf{0.0029} & 0.403 \\
\midrule
\multirow{4}{*}{AVIRIS} & 2D Spatial VAE & \textbf{0.0002} & 0.053 \\
 & 1D Pixelwise VAE & 0.0088 & 1.264 \\
 & 3D Spatio-Spectral VAE & 0.0047 & 1.540 \\
 & \textbf{PRISM} & 0.0084 & 1.334 \\
\midrule
\multirow{4}{*}{CRISM} & 2D Spatial VAE & 0.5697 & 15.134 \\
 & 1D Pixelwise VAE & 1.4521 & 13.854 \\
 & 3D Spatio-Spectral VAE & \textbf{0.3728} & 15.394 \\
 & \textbf{PRISM} & 0.5093 & 15.143 \\
\bottomrule

\end{tabular}
\caption{Chemical-interpolation smoothness: linearly interpolating between two patches' latent codes and decoding, jaggedness is the mean second-difference magnitude of the decoded spectrum along the interpolation path (lower = smoother); path length is the total variation of the decoded spectrum along the same path, reported for context rather than as a quality axis.}
\label{tab:table3}
\end{table}

Table~\ref{tab:table3} shows PRISM producing the smoothest interpolation path on IIRS by a wide margin and remaining competitive on AVIRIS, while the 2D Spatial VAE's very low jaggedness on AVIRIS is again a collapse artifact (a nearly-constant decoder trivially interpolates smoothly because it barely moves at all) rather than a genuine chemical-interpolation result -- the same caveat as Table~\ref{tab:table2}. On CRISM, where SAM's near-zero-energy floor is known to affect roughly half of all pixels (Appendix~\ref{app:diagnostics}), jaggedness is noisy across every model and no architecture shows a clear advantage. Taken together with the noise-injection result, chemical interpolation is where PRISM's smooth, low-dimensional, per-pixel spectral latent shows its intended benefit most clearly on the dataset with the least confounding degeneracy.

\begin{table}[t]
\centering
\small
\resizebox{\linewidth}{!}{%
\begin{tabular}{llcccccc}
\toprule
\textbf{Dataset} & \textbf{Model} & \textbf{Masked bands} & \textbf{PSNR (masked)$\uparrow$} & \textbf{Mean-fill PSNR} & \textbf{SAM-valid ($^\circ$)$\downarrow$} & \textbf{Rel.\ gain$\uparrow$} & \textbf{Pass-through} \\
\midrule
\input{tables/table4_inpainting}
\end{tabular}%
}
\caption{Spectral band-masking recovery: a contiguous block covering 10\% of a sensor's bands is zeroed at the model's input, at five positions spanning the spectrum, and reconstruction quality is scored only on the masked bands, against filling the same bands with the band-wise training-set mean. Deliberately a \emph{band}, not a pixel, mask: every architecture, including the 1D Pixelwise VAE (which has no spatial context to draw on at all), still has the pixel's other bands available, so this is the one masking scheme every architecture in the grid can be meaningfully scored on. None of these models was trained on masked input, so this measures an implicit spectral prior under distribution shift, not a trained inpainting capability. Pass-through is the model's masked-region MSE divided by a zero-fill null's masked-region MSE: a value near 1 indicates the model largely reproduced the zeros it was given at its input rather than drawing on a learned spectral prior, while a value well below 1 indicates it filled the gap from something it inferred about spectral shape. Cells marked ``\tbd{}'' await the pending re-evaluation pass noted in Table~\ref{tab:table1}.}
\label{tab:table4}
\end{table}

Table~\ref{tab:table4} is populated once the band-masking diagnostic introduced alongside this draft has been run against every checkpoint (\S\ref{sec:results}, mechanism diagnostics); we report the protocol here and will fill the cells in the accompanying revision. The comparison this table is designed to make is direct: an architecture whose reconstruction on the masked bands only beats a static band-wise mean fill by a wide margin has learned something about spectral covariance across bands that a model whose gain over that mean fill is small or negative has not. We expect, and will report plainly if disconfirmed, that PRISM's isolated spectral stream -- which never mixes pixels and is trained under the same SAM-bounded objective throughout -- retains more of this implicit spectral prior under masking than a model whose spectral information has been blurred across pixels by 2D or 3D spatial convolutions.

\paragraph{Combined read.} Across the four tables, PRISM's advantage is concentrated exactly where its architecture predicts it should be: it wins reconstruction fidelity on both spatial (PSNR) and spectral (SAM) axes simultaneously on every dataset (Table~\ref{tab:table1}), at a fraction of the parameter count of the largest baseline, without any single competing architecture matching it on both axes at once. Its one clear weakness -- raw sensitivity to latent-space noise (Table~\ref{tab:table2}) -- is a genuine property worth flagging for the downstream latent-diffusion-backbone use case the paper motivates, not a confound to explain away, and we report it as such rather than omit it. Chemical interpolation (Table~\ref{tab:table3}) is smoothest for PRISM on the dataset least affected by other architectures' collapse artifacts, consistent with the same low-dimensional, isolated spectral latent that gives PRISM its SAM advantage in Table~\ref{tab:table1} also giving it a smoother generative manifold. The overall picture supports the paper's central claim -- that structurally isolating spatial and spectral information, and re-coupling them only under a physics-bounded fusion, buys simultaneous spatial and spectral fidelity that no single-stream baseline achieves -- while being explicit about the one axis (raw noise robustness) where that same fidelity has a measurable cost.

\section{Conclusion}

We introduced PRISM, a dual-stream physics-informed VAE for hyperspectral imagery that structurally isolates spatial and spectral feature extraction and re-couples them only at a learned, SAM-bounded fusion. Across three planetary and terrestrial hyperspectral datasets and three baseline architectures spanning the 2D, pixelwise-1D, and 3D spatio-spectral design space, PRISM achieves the best or near-best reconstruction fidelity on both spatial (PSNR, SSIM) and spectral (SAM) axes simultaneously, at a fraction of the parameter count of the largest baseline, at a matched latent-rate bottleneck across every architecture -- a combination no single-stream baseline in our grid attains. Its decoupled, per-pixel spectral latent supports smooth chemical interpolation, motivating its use as a candidate backbone for a Latent Diffusion Model performing purification of deep-space transmission corruptions and generation of omitted spatial frames, though we are explicit that PRISM's high-fidelity decoder is measurably more sensitive to raw latent-space noise than some lower-fidelity baselines -- a genuine property to design around in that downstream use, not a result we obscure. The architecture, objective, and evaluation protocol carry over unchanged to any other physics-heavy domain where spatial denoising and per-pixel chemical or physical fidelity are both required, such as precision agriculture and climate and atmospheric monitoring.

\subsection*{AI use statement}

Yes, to aid or polish writing. Details are described in the paper. Yes, to draft sections of the paper. Details are described in the paper. Generative AI assistance was used to help draft and revise sections of this manuscript, including restructuring related-work and methodology prose, drafting table captions and discussion paragraphs from the authors' own experimental results and code, and copy-editing. All AI-drafted content was reviewed, verified against the underlying code and result files, and edited by the authors, who take full responsibility for the accuracy of every technical claim, number, and citation in this paper. Generative AI was not used to generate experimental results, run models, or fabricate data; every number reported in this paper is read from the authors' own result files.

\subsection*{Ethics statement}

Planetary hyperspectral imagery informs long-horizon scientific claims such as the presence or absence of water, minerals, and organic markers on other worlds. Any generative model that reconstructs or inpaints such data has some propensity to hallucinate: to place a mineral where there is none, or to remove one that is present. PRISM's SAM prior is designed to make such propensities costly, and our downstream-readiness benchmarks are designed to make degradation modes measurable. Nevertheless, we recommend that reconstructions produced by PRISM be used to \emph{complement} raw sensor data in scientific pipelines, not to replace it, and that mineralogical conclusions drawn from reconstructed cubes be flagged as such in downstream publications. The same considerations apply to any transfer of these methods to other domains, where high-stakes policy and resource decisions may be downstream of the reconstruction.

\subsection*{Reproducibility statement}

All five architectures compared in this paper (three baselines plus PRISM) implement one shared model-agnostic contract -- \texttt{forward}, \texttt{loss\_terms}, \texttt{reconstruct}, \texttt{encode\_latents}, \texttt{decode\_latents} -- registered in a single model registry, so that the training and evaluation code is identical across every cell of the ablation grid and never branches on model identity. Every hyperparameter reported in \S\ref{sec:objective}--\ref{sec:datasets} ($\beta$, $\lambda_{\mathrm{phys}}$, learning rate, batch size, epoch and patience budgets, and every model's latent-rate and reference-width knobs) is set in one per-dataset configuration file per sensor, rather than scattered across scripts, and the latent-rate matching described in \S\ref{sec:datasets} is verified programmatically against the actual instantiated models rather than only asserted. The falsification and downstream-readiness diagnostics reported in \S\ref{sec:results} and Appendix~\ref{app:diagnostics} read every threshold from a single preregistered configuration file rather than from inline constants, and every table in this paper is generated programmatically from the same underlying per-cell result files rather than transcribed by hand, so that a re-evaluation run reproduces every number in this paper without a manual editing step. All reconstruction, noise-injection, interpolation, and band-masking results are averaged over two independent training seeds (67 and 69) for every architecture and dataset reported.

\bibliography{../../references}
\bibliographystyle{iclr2027_conference}

\appendix
\section{Appendix}

\subsection{Mechanism diagnostics in full}
\label{app:diagnostics}

Table~\ref{tab:table5} reports the trivial-predictor floor diagnostic introduced in \S\ref{sec:results} in full: for every cell, the best zero-rate floor (the best of the train-set global, fold, and region mean spectra -- predictors that carry no per-patch information at all), the model's lift over that floor in PSNR and relative SAM-valid, its lift over the substantially stronger (but still learning-free) per-patch spatial-mean predictor, and the fraction of the gap between that floor and a perfect-copy identity oracle that the model's PSNR actually captures. As with the band-masking diagnostic in Table~\ref{tab:table4}, this is reported as a diagnostic number, with no pass/fail threshold: a previous version of this check, gating on a fixed minimum lift, flagged every trained cell in an earlier grid as invalid, because the per-patch spatial-mean predictor is a substantially stronger baseline than a global floor on smooth, highly spatially redundant hyperspectral data, and a $64{:}1$-compressed model losing to it does not by itself indicate a fault in the architecture.

\begin{table}[h]
\centering
\small
\resizebox{\linewidth}{!}{%
\begin{tabular}{llccccc}
\toprule
\textbf{Dataset} & \textbf{Model} & \textbf{Floor PSNR} & \textbf{Lift (dB)$\uparrow$} & \textbf{Lift SAM-v (rel.)$\uparrow$} & \textbf{vs.\ mean-patch (dB)} & \textbf{Headroom$\uparrow$} \\
\midrule
% GENERATED by paper/iclr/v2-claude/make_tables.py — DO NOT EDIT BY HAND
% source: /home/minikazuki/Documents/Research Papers/A Dual Stream Physics Informed VAE for hallucination free LDM purification of HSI/results/probes.csv (P1 trivial floors, diagnostic only) | seeds=[67, 69] loss=physics select=sam
\multirow{4}{*}{IIRS} & 2D Spatial VAE & \tbd{} & \tbd{} & \tbd{} & \tbd{} & \tbd{} \\
 & 1D Pixelwise VAE & \tbd{} & \tbd{} & \tbd{} & \tbd{} & \tbd{} \\
 & 3D Spatio-Spectral VAE & \tbd{} & \tbd{} & \tbd{} & \tbd{} & \tbd{} \\
 & \textbf{PRISM} & \tbd{} & \tbd{} & \tbd{} & \tbd{} & \tbd{} \\
\midrule
\multirow{4}{*}{AVIRIS} & 2D Spatial VAE & \tbd{} & \tbd{} & \tbd{} & \tbd{} & \tbd{} \\
 & 1D Pixelwise VAE & \tbd{} & \tbd{} & \tbd{} & \tbd{} & \tbd{} \\
 & 3D Spatio-Spectral VAE & \tbd{} & \tbd{} & \tbd{} & \tbd{} & \tbd{} \\
 & \textbf{PRISM} & \tbd{} & \tbd{} & \tbd{} & \tbd{} & \tbd{} \\
\midrule
\multirow{4}{*}{CRISM} & 2D Spatial VAE & \tbd{} & \tbd{} & \tbd{} & \tbd{} & \tbd{} \\
 & 1D Pixelwise VAE & \tbd{} & \tbd{} & \tbd{} & \tbd{} & \tbd{} \\
 & 3D Spatio-Spectral VAE & \tbd{} & \tbd{} & \tbd{} & \tbd{} & \tbd{} \\
 & \textbf{PRISM} & \tbd{} & \tbd{} & \tbd{} & \tbd{} & \tbd{} \\
\bottomrule

\end{tabular}%
}
\caption{Trivial-predictor floor diagnostic (P1), diagnostic only, no pass/fail threshold. Cells marked ``\tbd{}'' await the pending re-evaluation pass noted in the main text.}
\label{tab:table5}
\end{table}

The posterior-collapse and spatial-reliance-shuffle diagnostics summarized in \S\ref{sec:results} identify the 2D Spatial VAE under the standard-ELBO objective as the architecture most prone to collapse on IIRS and AVIRIS in this grid -- consistent with that architecture having no physics term to prevent a decoder from settling on a near-constant, low-error output when the reconstruction objective alone permits it -- which is the collapse referenced when reading Table~\ref{tab:table2} and Table~\ref{tab:table3}'s negative or anomalously low values for that model. The 1D Pixelwise VAE's spatial-reliance-shuffle score, by construction, differs from its intact score by a numerically negligible amount (within the diagnostic's own preregistered tolerance for exact permutation-equivariance), confirming the diagnostic itself is measuring what it claims to before any other model's shuffle result is read as a finding.

\subsection{Reconstruction metrics by seed}
Per-seed (rather than seed-averaged) reconstruction metrics, and the corresponding pairwise bootstrap/permutation/Holm-corrected significance tests, are available in the accompanying result files and will be included in tabular form in a revision.

\end{document}
